\chapter*{คำนำ}
\addcontentsline{toc}{chapter}{คำนำ}

เอกสารคู่มือการประกอบ ติดตั้ง และใช้งานเครื่องสเปกโทรโฟโตมิเตอร์ตรวจวัดธาตุอาหารหลักในดินและสมบัติเชิงแสงของของเหลว (NPK Level Meter 1.02 \& Liquid Optics Analyzer) เล่มนี้ จัดทำขึ้นเพื่อเป็นคู่มือวิชาการและแนวทางการปฏิบัติการเชิงลึก สำหรับนักวิจัย คณาจารย์ นักศึกษา และเกษตรกรผู้สนใจนวัตกรรมเกษตรแม่นยำ (Precision Agriculture) โดยเฉพาะการบริหารจัดการธาตุอาหารในสวนทุเรียนแปลงใหญ่ในเขตภาคตะวันออกของประเทศไทย

ระบบตรวจวัดนี้ได้รับการออกแบบและพัฒนาขึ้นภายใต้แนวคิดการประมวลผลบนชิปสมบูรณ์ (Edge-Only Architecture) โดยผสานสถาปัตยกรรมไมโครคอนโทรลเลอร์สมรรถนะสูง Microchip ATSAMD51 ARM Cortex-M4F ร่วมกับเซนเซอร์วัดค่าสีดิจิทัล TCS34725 และแหล่งกำเนิดแสงสเปกตรัมหลายช่วงคลื่น NeoPixel WS2812B ทำให้สามารถวิเคราะห์ธาตุอาหารดิน ไนโตรเจน (N) ฟอสฟอรัส (P) และโพแทสเซียม (K) ได้อย่างรวดเร็ว พร้อมทั้งรองรับการคำนวณ 3 โหมด คือ โครงข่ายประสาทเทียม TinyML Deep Learning, สมการพหุนามสากล, และระบบสร้างเส้นโค้งมาตรฐานสดในตัวเครื่อง (In-Situ Standard Curve Wizard) ที่คำนวณความชัน จุดตัด $R^2$ LOD และ LOQ ตามเกณฑ์สากลของ IUPAC ตลอดจนการขยายขีดความสามารถสู่หน้าจอที่ 5 เพื่อตรวจวัดดัชนีหักเหของแสง ($n$) ความหนาแน่น ($\rho$) ปริมาณของแข็ง ($^\circ\text{Bx}$) และระดับความขุ่นของตัวอย่างของเหลวทางการเกษตร

เนื้อหาภายในคู่มือเล่มนี้แบ่งออกเป็น 5 บทหลัก และ 2 ภาคผนวก ประกอบด้วย
\begin{enumerate}[leftmargin=1.5em, itemsep=3pt]
    \item \textbf{บทที่ 1 บทนำและหลักการฟิสิกส์เชิงแสง} อธิบายทฤษฎีกฎของเบียร์และแลมเบิร์ต สเปกตรัมการดูดกลืนแสง และแบบจำลองมาตรวิทยาของเหลว
    \item \textbf{บทที่ 2 การออกแบบและประกอบชิ้นส่วนฮาร์ดแวร์} แสดงรายการอุปกรณ์ แผนผังการเชื่อมต่อวงจร และขั้นตอนการประกอบเชิงกล
    \item \textbf{บทที่ 3 การติดตั้งเฟิร์มแวร์และการตั้งค่าซอฟต์แวร์} แนะนำไลบรารีที่จำเป็น การตั้งค่าฟอนต์ภาษาไทย และการคอมไพล์ผ่าน Arduino CLI
    \item \textbf{บทที่ 4 คู่มือการปฏิบัติการและแบบจำลองการเทียบมาตรฐาน} นำเสนอการทำงานของระบบ 5 หน้าจอ ขั้นตอนการวัดดินและของเหลว และสถิติการสร้างกราฟมาตรฐาน
    \item \textbf{บทที่ 5 การบำรุงรักษาและการแก้ไขปัญหา} รวบรวมแนวทางการดูแลรักษาอุปกรณ์และการแก้ไขข้อขัดข้องที่อาจเกิดขึ้นในภาคสนาม
    \item \textbf{ภาคผนวก} รวบรวมตารางค่าดัชนีหักเหและความหนาแน่นมาตรฐาน ICUMSA และข้อมูลจำเพาะเชิงสถาปัตยกรรมของ Wio Terminal
\end{enumerate}

คณะผู้จัดทำหวังเป็นอย่างยิ่งว่า คู่มือวิชาการเล่มนี้จะเป็นประโยชน์ต่อการนำไปประยุกต์ใช้งานจริงในภาคสนาม การจัดการเรียนการสอนในห้องปฏิบัติการฟิสิกส์และเครื่องมือวัด และเป็นรากฐานในการพัฒนาเทคโนโลยีเกษตรอัจฉริยะของประเทศไทยสืบไป

\vspace{1.0cm}
\begin{flushright}
    \begin{minipage}{0.55\textwidth}
        \centering
        \textbf{ผู้ช่วยศาสตราจารย์ ดร.ชีวะ ทัศนา และคณะ}\\
        หน่วยวิจัยฟิสิกส์เพื่อการเกษตรดิจิทัลและปัญญาประดิษฐ์ (AI4D AgriPhysics)\\
        สาขาวิชาฟิสิกส์ คณะวิทยาศาสตร์และเทคโนโลยี\\
        มหาวิทยาลัยราชภัฏรำไพพรรณี\\
        กันยายน 2569
    \end{minipage}
\end{flushright}
